\chapter{Eksponensiale}
\setcounter{figure}{1}
\setcounter{subfigure}{1}


$ \hspace{-5pt}\begin{array}{cccccccccccc}   \end{array} $ \hspace{2 pt}\raisebox{-0.2em}{\includegraphics[height=1em]{../icons/www.pdf}} {(section shortcode: MG10042 )} \par 



\label{m38359*cid3} $ \hspace{-5pt}\begin{array}{cccccccccccc}   \includegraphics[width=0.75cm]{col11306.imgs/summary_video.png} &   \end{array} $ \hspace{2 pt}\raisebox{-5 pt}{} {(section shortcode: MG10043 )} \par 
\label{m38359*id62562}Eksponensiaalnotasie is ’n kort manier om te skryf dat ’n getal meermale met homself vermenigvuldig word. Laat ons beter definieer hoe om eksponensiaalnotasie te gebruik

 \begin{Large}
\begin{center}
$ _{\mbox{grondtal}~\leftarrow} $\begin{Large} $ ~a^{n~\rightarrow~$ \end{Large}$\mbox{eksponent / index}} $
\end{center}
 \end{Large}

For any real number $a$ and natural number $n$, we can write $a$ multiplied by itself $n$ times as $a^n$.

 
\Identity{Definisie:}{
\begin{flushleft}
\begin{enumerate}[itemsep=5pt, label=\textbf{\arabic*}.]
 \item $a^n = a \times a \times a \times \ldots \times a ~~ (n ~ \mbox{keer}) ~~~~ (a \in \mathbb{R}, n \in \mathbb{N})$
 \item $ a^0 = 1 \hspace{1cm}$  ($a \ne 0 $ because $0^0$ is undefined)
 \item $a^{-n} &=& \frac{1}{a^n}\hspace{0.5cm}$ ($a \ne 0 $ because $\frac{1}{0}$ is undefined)
\end{enumerate}



\end{flushleft}

} 

Voorbeelde:
\begin{enumerate}[noitemsep, label=\textbf{\arabic*.}]
\item $3 \times 3 = 3^2$
\item $5 \times 5 \times 5 \times 5 = 5^4 $
\item $p \times p \times p = p^3$
\item $(3^x)^0 = 1$
\item $ 2^{-4} = \dfrac{1}{2^4} = \frac{1}{16}$
\item $ \dfrac{1}{5^{-x}} = 5^x$
\end{enumerate}

% We now expand the definition to include 0 and the negative integers.
% 
% If exponent $n$ is 0, then we have 
% 
% 
% 
% \Tip{The restriction applies because $0^0$ is undefined }
% 
% We also define what it means if we have a negative exponent $a^{-n}$.
% \begin{eqnarray*}
%     a^{-n} &=& \frac{1}{a^n} \\
%            &=& \frac{1}{a \times a \times a \times \ldots \times a ~~ (n ~ \textrm{times})} 
% \end{eqnarray*}
% 
% 
% \Tip{Here $a$ cannot equal 0 because $\frac{1}{0}$ is undefined}      
      
\textbf{Note:} We always write the final answer with positive exponents.


% If exponent $n$ is an even integer, then ${a}^{n}$ will always be positive for any non-zero real number $a$. 
% 
% For example, although $-3$ is negative, but $(-3)^2=-3 \times -3 = 9$ which is positive and $(-3)^{-2} = \frac{1}{-3 \times -3} = \frac{1}{9} $ is also positive.
% 
% If the exponent $n$ is an odd integer, then for any non-zero real number $a$ 
% 
% \begin{eqnarray*}
% a^n ~~ \mbox{is positive if} ~~ a > 0 \\
% a^n ~~ \mbox{is negative if} ~~ a < 0
% \end{eqnarray*}
% 
% For example, $(-2)^3 = -2 \times -2 \times -2 = -8$ and $(2)^{-5} = \dfrac{1}{2 \times 2 \times 2 \times 2 \times 2} = \dfrac{1}{32}$.

\setcounter{subfigure}{0}
\begin{figure}[H] % horizontal\label{m38359*Exponents-1}
\textnormal{Khan Academy video on Exponents - 1}\vspace{.1in} 
\label{m38359*yt-media1}\label{m38359*yt-video1}
\raisebox{-5 pt}{ \includegraphics[width=0.5cm]{col11306.imgs/summary_www.png}} { (Video:  MG10044 )}
% \vspace{2pt}
% \vspace{.1in}
\end{figure}       


\setcounter{subfigure}{0}
\begin{figure}[H] % horizontal\label{m38359*Exponents-2}
\textnormal{Khan Academy video on Exponents-2}\vspace{.1in} 
\label{m38359*yt-media2}\label{m38359*yt-video2}
\raisebox{-5 pt}{ \includegraphics[width=0.5cm]{col11306.imgs/summary_www.png}} { (Video:  MG10045 )}
% \vspace{2pt}
% \vspace{.1in}
\end{figure}       




\section { Laws of exponents}
\nopagebreak
\label{m38359*cid4} $ \hspace{-5pt}\begin{array}{cccccccccccc}   \includegraphics[width=0.75cm]{col11306.imgs/summary_fullmarks.png} &   \includegraphics[width=0.75cm]{col11306.imgs/summary_video.png} &   \end{array} $ \hspace{2 pt}\raisebox{-5 pt}{} {(section shortcode: MG10046 )} \par 
\label{m38359*id63061}Daar is heelwat eksponentwette wat ons kan gebruik om getalle met eksponente te vereenvoudig. Sommige
van hierdie wette het ons reeds in vorige grade teëgekom, maar ons sal die volledige lys hier sien en elke wet
verduidelik, sodat jy hulle kan verstaan en nie bloot memoriseer nie


\begin{eqnarray*}
{a}^{m}\ensuremath{\times}{a}^{n}& =& {a}^{m+n}\hfill \\ 
\frac{{a}^{m}}{{a}^{n}}& =& {a}^{m-n}\hfill \\ 
{\left(ab\right)}^{n}& =& {a}^{n}{b}^{n}\hfill \\ 
        (\frac{a}{b})^n & = & \frac{a^n}{b^n} \\
{\left({a}^{m}\right)}^{n}& =& {a}^{mn}\hfill 
\end{eqnarray*}
waar $a > 0, b > 0$ en $m,n \in \mathbb{Z}$

% 
%       \label{m38359*uid4}
%             \  \subsubsubsection { Exponential Law 1: ${a}^{0}=1$}
%             \nopagebreak
%         \label{m38359*id63512}Our definition of exponential notation shows that\par 
%         \label{m38359*uid5}\nopagebreak\noindent{}
%           
%     \begin{center}
%     \begin{array}{ccl}\hfill {a}^{0}& =& 1,\left(a\ne 0\right)\hfill \end{array}
%       \end{center}
%         \label{m38359*eip-662}To convince yourself of why this is true, use the fourth exponential law above (division of exponents) and consider what happens when $m=n$.\par \label{m38359*id63571}For example, ${x}^{0}=1$ and ${\left(1\phantom{\rule{0.277778em}{0ex}}000\phantom{\rule{0.277778em}{0ex}}000\right)}^{0}=1$.\par 
% \label{m38359*secfhsst!!!underscore!!!id339}
%               \subsubsubsection { Exercise: Exponential Law 1: ${a}^{0}=1,\left(a\ne 0\right)$ }
%             \nopagebreak
%         \label{m38359*id63666}\begin{enumerate}[noitemsep, label=\textbf{\arabic*}. ] 
%             \label{m38359*uid6}\item 
%             ${16}^{0}$
%       \label{m38359*uid7}\item 
%         $16{a}^{0}$
%       \label{m38359*uid8}\item 
%         ${\left(16+a\right)}^{0}$
%       \label{m38359*uid9}\item 
%         ${\left(-16\right)}^{0}$
%       \label{m38359*uid10}\item 
%         $-{16}^{0}$ 
% \newline
% \newline
%           \end{enumerate}
%       \label{m38359*uid11}
% \par \raisebox{-0.2em}{\includegraphics[height=1em]{../icons/www.pdf}} Find the answers with the shortcodes:
%  \par \begin{tabular}[h]{cccccc}
%  (1.) lOG  & \end{tabular}
%             \subsubsubsection{ Exponential Law 2: 

% \begin{center}
%   \begin{array}[lc]
%       \textrm{\textbf{Exponential Law:}} & c
%   \end{array}
% \end{center}

\Identity{Exponential Law}
{
$${a}^{m}\ensuremath{\times}{a}^{n}={a}^{m+n}$$

Die definisie van eksponensiaalnotasie wys dat

\begin{eqnarray*}
	      {a}^{m}\ensuremath{\times}{a}^{n}& =& a\ensuremath{\times}a\ensuremath{\times} \ldots \ensuremath{\times}a\hfill & \left(\mbox{$m$ keer}\right)\hfill \\ 
	      & & \phantom{\rule{-0.166667em}{0ex}}\phantom{\rule{-0.166667em}{0ex}}\phantom{\rule{-0.166667em}{0ex}}\phantom{\rule{-0.166667em}{0ex}}\ensuremath{\times}a\ensuremath{\times}a\ensuremath{\times}\ldots\ensuremath{\times}a\hfill & \left(\mbox{$n$ keer}\right)\hfill \\ 
	      & =& a\ensuremath{\times}a\ensuremath{\times}\ldots\ensuremath{\times}a\hfill & \left(\mathrm{m}+\mbox{$n$ keer}\right)\hfill \\ 
	      & =& {a}^{m+n}\hfill & 
\end{eqnarray*}



}


\label{m38359*id64082}Byvoorbeeld,\par 


\begin{eqnarray*}
{2}^{4}\ensuremath{\times}{2}^{3}& =& \left(2\ensuremath{\times}2\ensuremath{\times}2\ensuremath{\times}2\right)\ensuremath{\times}\left(2\ensuremath{\times}2\ensuremath{\times}2\right)\hfill \\
	      & =& {2}^{4+3} \\
	      & =& {2}^{7}
\end{eqnarray*}

\Note{Hierdie eenvoudige wet is die rede waarom eksponente oorspronklik geskep
is. Voor die dae van rekenaars moes vermenigvuldiging met potlood en papier gedoen
word. Dit vat baie lank om vermenigvuldiging te doen, maar dit is vinnig en eenvoudig
om getalle bymekaar te tel. Hierdie eksponentwet wys dat dit moontlik is om twee
getalle te vermenigvuldig deur hulle eksponente bymekaar te tel (indien hulle dieselfde
grondtal het). Hierdie ontdekking het wiskundiges baie tyd gespaar, wat hulle toe kon
gebruik om iets meer produktiefs te doen.}

\begin{figure}[H] % horizontal\label{m38359*ExponentsRule1}
\textnormal{Khan Academy video on Exponents - 3}\vspace{.1in} \nopagebreak
\label{m38359*yt-media3}\label{m38359*yt-video3}
\raisebox{-5 pt}{ \includegraphics[width=0.5cm]{col11306.imgs/summary_www.png}} { (Video:  MG10047 )}
% \vspace{2pt}
% \vspace{.1in}
\end{figure}
    
% \par
% \subsubsubsection{  Application using Exponential Law 2: ${a}^{m}\ensuremath{\times}{a}^{n}={a}^{m+n}$ }
%             \nopagebreak
%         \label{m38359*id64269}\begin{enumerate}[noitemsep, label=\textbf{\arabic*}. ] 
%             \label{m38359*uid13}\item 
%             ${x}^{2}\ensuremath{\cdot}{x}^{5}$
%       \label{m38359*uid14}\item 
%         ${2}^{3}\ensuremath{\cdot}{2}^{4}$
%         [Take note that the base (2) stays the same.]
%       \label{m38359*uid15}\item 
%         $3\ensuremath{\times}{3}^{2a}\ensuremath{\times}{3}^{2}$
% \newline
% \newline
%           \end{enumerate}
%       \label{m38359*uid16}
% \par \raisebox{-0.2em}{\includegraphics[height=1em]{../icons/www.pdf}} Find the answers with the shortcodes:
%  \par \begin{tabular}[h]{cccccc}
%  (1.) lO7  & \end{tabular}
%             \subsubsubsection{ Exponential Law 3: ${a}^{-n}=\frac{1}{{a}^{n}},a\ne 0$}
%             \nopagebreak
%         \label{m38359*id64482}Our definition of exponential notation for a negative exponent shows that\par 
%         \label{m38359*uid17}\nopagebreak\noindent{}
%     \begin{center}
%     \begin{array}{cccc}\hfill {a}^{-n}& =& 1÷a÷...÷a\hfill & \left(\mathrm{n\; times}\right)\hfill \\ \hfill & =& \frac{1}{1\ensuremath{\times}a\ensuremath{\times}\cdots \ensuremath{\times}a}\hfill & \left(\mathrm{n\; times}\right)\hfill \\ \hfill & =& \frac{1}{{a}^{n}}\hfill & \end{array}\tag{5.8}
%       \end{center}
%         \label{m38359*id64624}This means that a minus sign in the exponent is just another way of showing that the whole exponential number is to be divided instead of multiplied.\par 
%         \label{m38359*id64630}For example,\par 
%         \label{m38359*id64634}\nopagebreak\noindent{}
%           
%     \begin{center}
%     \begin{array}{ccl}\hfill {2}^{-7}& =& \frac{1}{2\ensuremath{\times}2\ensuremath{\times}2\ensuremath{\times}2\ensuremath{\times}2\ensuremath{\times}2\ensuremath{\times}2}\hfill \\ & =& \frac{1}{{2}^{7}}\hfill \end{array}\tag{5.9}
%       \end{center}
% \label{m38359*eip-294}This law is useful in helping us simplify fractions where there are exponents in both the denominator and the numerator. For example:
% \label{m38359*id7902}\nopagebreak\noindent{}
% 
%     \begin{center}
%     \begin{array}{ccl}\frac{{a}^{-3}}{{a}^{4}}& =& \frac{1}{{a}^{3}{a}^{4}}\\ & =& \frac{1}{{a}^{7}}\end{array}\tag{5.10}
%       \end{center}
% \par \label{m38359*secfhsst!!!underscore!!!id835}
%             \subsubsubsection{  Application using Exponential Law 3: ${a}^{-n}=\frac{1}{{a}^{n}},a\ne 0$ }
%             \nopagebreak
%         \label{m38359*id64771}\begin{enumerate}[noitemsep, label=\textbf{\arabic*}. ] 
%             \label{m38359*uid18}\item 
%             ${2}^{-2}=\frac{1}{{2}^{2}}$
%       \label{m38359*uid19}\item 
%         $\frac{{2}^{-2}}{{3}^{2}}$
%       \label{m38359*uid20}\item 
%         ${\left(\frac{2}{3}\right)}^{-3}$
%       \label{m38359*uid21}\item 
%         $\frac{m}{{n}^{-4}}$
%       \label{m38359*uid22}\item 
%         $\frac{{a}^{-3}\ensuremath{\cdot}{x}^{4}}{{a}^{5}\ensuremath{\cdot}{x}^{-2}}$
% \newline
% \newline
%           \end{enumerate}
%       \label{m38359*uid23}
% \par \raisebox{-0.2em}{\includegraphics[height=1em]{../icons/www.pdf}} Find the answers with the shortcodes:
%  \par \begin{tabular}[h]{cccccc}
%  (1.) lcx  & \end{tabular}

% \begin{center}
% \begin{array}[lc]
% \textrm{\textbf{Exponential Law:}} & $$ \frac{ {a}^{m} }{ {a}^{n} }={a}^{m-n}$$
% \end{array}
% \end{center}


\Identity{Eksponente wet}
{
$$ \frac{ {a}^{m} }{ {a}^{n} }={a}^{m-n}$$
\begin{eqnarray*}
\frac{a^m}{a^n} &=& \frac{a \times a \times a \ldots \times a ~ (m~\mbox{keer})} {a \times a \times a \ldots \times a ~ (n~\mbox{keer})} \\
&= & a \times a \times a \ldots \times a ~ (m-n~\mbox{keer}) \\
&= & a^{m-n}
\end{eqnarray*}

}




          
\label{m38359*id65293}For example,

\begin{eqnarray*}
    \dfrac{{2}^{7}}{{2}^{3}}& =& \dfrac{2\ensuremath{\times}2\ensuremath{\times}2\ensuremath{\times}2\ensuremath{\times}2\ensuremath{\times}2\ensuremath{\times}2}{2\ensuremath{\times}2\ensuremath{\times}2}\hfill \\
						     & =& 2\ensuremath{\times}2\ensuremath{\times}2\ensuremath{\times}2\hfill \\
						     & =& {2}^{4}\hfill \hfill 
\end{eqnarray*}

\setcounter{subfigure}{0}
\begin{figure}[H] % horizontal\label{m38359*exponents-5}
\textnormal{Khan academy video on exponents - 4}\vspace{.1in} \nopagebreak
\label{m38359*yt-media6}\label{m38359*yt-video6}
\raisebox{-5 pt}{ \includegraphics[width=0.5cm]{col11306.imgs/summary_www.png}} { (Video:  MG10048 )}
% \vspace{2pt}
% \vspace{.1in}
\end{figure}       
% \par \label{m38359*secfhsst!!!underscore!!!id1192}
%             \subsubsubsection{ Application using Exponential Law 4: ${a}^{m}÷{a}^{n}={a}^{m-n}$}
%             \nopagebreak
%         \label{m38359*id65493}\begin{enumerate}[noitemsep, label=\textbf{\arabic*}. ] 
%             \label{m38359*uid25}\item 
%             $\frac{{a}^{6}}{{a}^{2}}={a}^{6-2}$
%       \label{m38359*uid26}\item 
%         $\frac{{3}^{2}}{{3}^{6}}$
%       \label{m38359*uid27}\item 
%         $\frac{32{a}^{2}}{4{a}^{8}}$
%       \label{m38359*uid28}\item 
%         $\frac{{a}^{3x}}{{a}^{4}}$
% \newline8
% \newline
%           \end{enumerate}
%       \label{m38359*uid29}
% \par \raisebox{-0.2em}{\includegraphics[height=1em]{../icons/www.pdf}} Find the answers with the shortcodes:
%  \par \begin{tabular}[h]{cccccc}
%  (1.) lOA  & \end{tabular}


% \begin{center}
%   \begin{array}[lc]
%       \textrm{\textbf{Exponential Law:}} & $$ {\left(ab\right)}^{n}={a}^{n}{b}^{n}$$
%   \end{array}
% \end{center}

\Identity{Exponential Law}
{
 $$ {\left(ab\right)}^{n}={a}^{n}{b}^{n}$$

\par
Die volgorde waarin twee getalle vermenigvuldig word, is onbelangrik. Dus $a \times b = b \times a$ 
\begin{center}
\begin{array}{cclc}\hfill {\left(ab\right)}^{n}& =& a\ensuremath{\times}b\ensuremath{\times}a\ensuremath{\times}b\ensuremath{\times}\ldots\ensuremath{\times}a\ensuremath{\times}b\hfill & \left(\mbox{$n$ times}\right)\hfill \\
	\hfill & =& a\ensuremath{\times}a\ensuremath{\times}\ldots\ensuremath{\times}a\hfill & \left(\mbox{$n$ keer}\right)\hfill \\
	\hfill & & \phantom{\rule{-0.166667em}{0ex}}\phantom{\rule{-0.166667em}{0ex}}\phantom{\rule{-0.166667em}{0ex}}\phantom{\rule{-0.166667em}{0ex}}\ensuremath{\times}b\ensuremath{\times}b\ensuremath{\times}\ldots\ensuremath{\times}b\hfill & \left(\mbox{$n$ keer}\right)\hfill \\
	\hfill & =& {a}^{n}{b}^{n}\hfill & 
\end{array}
\end{center}
}


%             \subsubsubsection{ Exponential Law 5: ${\left(ab\right)}^{n}={a}^{n}{b}^{n}$}


\label{m38359*id66030}Byvoorbeeld,

\begin{center}
    \begin{array}{ccl}\hfill {\left(2\ensuremath{\cdot}3\right)}^{4}& =& \left(2\ensuremath{\cdot}3\right)\ensuremath{\times}\left(2\ensuremath{\cdot}3\right)\ensuremath{\times}\left(2\ensuremath{\cdot}3\right)\ensuremath{\times}\left(2\ensuremath{\cdot}3\right)\hfill \\
	    & =& \left(2\ensuremath{\times}2\ensuremath{\times}2\ensuremath{\times}2\right)\ensuremath{\times}\left(3\ensuremath{\times}3\ensuremath{\times}3\ensuremath{\times}3\right)\hfill \\
	    & =& \left({2}^{4}\right)\ensuremath{\times}\left({3}^{4}\right)\hfill \\ & =& {2}^{4}\cdot{3}^{4}\hfill 
    \end{array}
\end{center}



% \begin{center}
% \begin{array}[lc]
%   \textrm{\textbf{Exponential Law:}} & $$ (\frac{a}{b})^n = \frac{a^n}{b^n} $$
% \end{array}
% \end{center}


\Identity{Exponential Law}
{
$$ \Big(\frac{a}{b}\Big)^n = \frac{a^n}{b^n} $$

\begin{eqnarray*}

 \Big(\frac{a}{b}\Big)^n & = & \Big(\frac{a}{b}\Big) \times \Big(\frac{a}{b}\Big) \times \Big(\frac{a}{b}\Big) \times \ldots \times \Big(\frac{a}{b}\Big) ~~~(n~\mbox{times}) \\
                         & = & \frac{a \times a \times a \times \ldots \times a ~~~(n~\mbox{times})}{b \times b \times b \times \ldots \times b ~~~(n~\mbox{times})}\\
                         & = & \frac{a^n}{b^n}
\end{eqnarray*}



}




Byvoorbeeld,

\begin{eqnarray*}

\Big(\frac{2}{3}\Big)^3 & = & \Big(\frac{2}{3}\Big) \times  \Big(\frac{2}{3}\Big) \times \Big(\frac{2}{3}\Big) \\
                        & = & \frac{2 \times 2 \times 2}{3 \times 3 \times 3} \\
		        & = & \frac{2^3}{3^3}




\end{eqnarray*}







%             \subsubsubsection{ Application using Exponential Law 5: ${\left(ab\right)}^{n}={a}^{n}{b}^{n}$ }
%             \nopagebreak
%         \label{m38359*id66288}\begin{enumerate}[noitemsep, label=\textbf{\arabic*}. ] 
%             \label{m38359*uid31}\item 
%             ${\left(2xy\right)}^{3}={2}^{3}{x}^{3}{y}^{3}$
%       \label{m38359*uid32}\item 
%         ${\left(\frac{7a}{b}\right)}^{2}$
%       \label{m38359*uid33}\item 
%         ${\left(5a\right)}^{3}$
% \newline
% \newline
%           \end{enumerate}
%       \label{m38359*uid34}
% \par \raisebox{-0.2em}{\includegraphics[height=1em]{../icons/www.pdf}} Find the answers with the shortcodes:
%  \par \begin{tabular}[h]{cccccc}
%  (1.) lOs  & \end{tabular}
            

% \begin{center}
% \begin{array}[lc]
%   \textrm{\textbf{Exponential Law:}} & $$ {\left({a}^{m}\right)}^{n}={a}^{mn} $$
% \end{array}
% \end{center}


\Identity{Exponential Law}
{
$$ {\left({a}^{m}\right)}^{n}={a}^{mn} $$

Dit is moontlik om die eksponensiaal van ’n eksponensiaal te bereken. So, selfs al klink die eerste sin ingewikkeld, beteken dit bloot dat ’n mens die eksponensiaal van ’n getal
bereken en dan die eksponensiaal van die resultaat bereken. \par


\begin{center}
    \begin{array}{ccll}\hfill {\left({a}^{m}\right)}^{n}& =& {a}^{m}\ensuremath{\times}{a}^{m}\ensuremath{\times}\ldots\ensuremath{\times}{a}^{m}\hfill & \left(\mbox{$n$ keer}\right)\hfill \\
	\hfill & =& a\ensuremath{\times}a\ensuremath{\times}\ldots\ensuremath{\times}a\hfill & \left({m}\ensuremath{\times}\mbox{$n$ keer}\right)\hfill \\
	\hfill & =& {a}^{mn}\hfill & 
    \end{array}
\end{center}

}

% \subsubsubsection{ Exponential Law 6: ${\left({a}^{m}\right)}^{n}={a}^{mn}$}
%             \nopagebreak

\label{m38359*id66694}Byvoorbeeld,

\begin{center}
    \begin{array}{ccl}\hfill {\left({5}^{2}\right)}^{3}& =& \left({5}^{2}\right)\ensuremath{\times}\left({5}^{2}\right)\ensuremath{\times}\left({5}^{2}\right)\hfill \\ 
	      & =& \left(5\ensuremath{\times}5\right)\ensuremath{\times}\left(5\ensuremath{\times}5\right)\ensuremath{\times}\left(5\ensuremath{\times}5\right)\hfill \\
	      & =& \left({5}^{6}\right)\hfill
    \end{array}
\end{center}

\Note{

\begin{center}
    5^2 \times 5^4 & = & 5^{2+4} & = & 5^6  \\
\end{center}

\begin{center}
    (5^2)^4        & = & 5^{2\times4} & = & 5^8 
\end{center}

}

      
% \label{m38359*secfhsst!!!underscore!!!id1894}
%             \subsubsubsection{ Application using Exponential Law 6: ${\left({a}^{m}\right)}^{n}={a}^{mn}$ }
%             \nopagebreak
%         \label{m38359*id66924}\begin{enumerate}[noitemsep, label=\textbf{\arabic*}. ] 
%             \label{m38359*uid36}\item 
%             ${\left({x}^{3}\right)}^{4}$
%       \label{m38359*uid37}\item 
%         ${\left[{\left({a}^{4}\right)}^{3}\right]}^{2}$
%       \label{m38359*uid38}\item 
%         ${\left({3}^{n+3}\right)}^{2}$
% \newline
% \newline
%           \end{enumerate}
% \label{m38359*eip-323}\par
%             \label{m38359*secfhsst!!!underscore!!!id41}\vspace{.5cm} 

\begin{wex}
{ %title
Toepassing van eksponentwette
}
{%question
Vereenvoudig:

\begin{enumerate}[noitemsep, label=\textbf{\arabic*}.]
\item  $2^{3x} \times 2^{4x}$
\\
 \item $\frac{12p^2t^5}{3pt^3}$
\\
 \item $ (3x)^2 $
\\
 \item $(3^4\cdot5^2)^3$\\
\end{enumerate}

}
{%answer

\begin{enumerate}[noitemsep, label=\textbf{\arabic*}.]
\item  $2^{3x} \times 2^{4x} = 2^{3x+4x} = 2^{7x}$
\\
 \item $\frac{12p^2t^5}{3pt^3} = 4p^{(2-1)}t^{(5-3)} = 4pt^2$
\\
 \item $ (3x)^2 = 3^2x^2 = 9x^2$
\\
 \item $(3^4\cdot5^2)^3 = 3^{(4\times3)}\cdot5^{(2\times3)} = 3^{12}\cdot5^6  $
\end{enumerate}



}


\end{wex}









\begin{wex}
{
Exponential expressions
}
{
Vereenvoudig: $$\frac{2^{2n} \cdot 4^n \cdot 2 }{ 16^n} $$
}
{

\westep{Change bases to prime numbers}

$$ \frac{2^{2n} \cdot 4^n \cdot 2 }{ 16^n} ~~ = ~~ \frac{2^{2n} \cdot (2^2)^n \cdot 2^1 }{ (2^4)^n} $$


\westep{Vereenvoudig die eksponente}

\begin{eqnarray*}
 &= & \frac{ 2^{2n} \cdot 2^{2n} \cdot 2^1 }{ 2^{4n} } \\
& = & \frac{ 2^{2n + 2n +1}}{2^{4n}} \\
& = & \frac{2^{4n+1}}{2^{4n}} \\
& = & 2^{4n+1-(4n)} \\
& = & 2
\end{eqnarray*}
 
}
\end{wex}


     
\begin{wex}
{%title
Exponential expressions
}
{
Vereenvoudig: $\dfrac{{5}^{2x-1}\ensuremath{\cdot}{9}^{x-2}}{{15}^{2x-3}}$
}
{
%     \begin{enumerate}[noitemsep, label=\textbf{Step} \textbf{\arabic*}. ] 
%             \leftskip=20pt\rightskip=\leftskip\item  
%     \label{m38359*id118023}\nopagebreak\noindent{}
\westep{Change bases to prime numbers}
      
\begin{center}
\begin{array}{lcl} \dfrac{{5}^{2x-1}\ensuremath{\cdot}{9}^{x-2}}{{15}^{2x-3}}& =& \dfrac{{5}^{2x-1}\ensuremath{\cdot}{\left({3}^{2}\right)}^{x-2}}{{\left(5.3\right)}^{2x-3}}\hfill \\
		  & =& \dfrac{{5}^{2x-1}\ensuremath{\cdot}{3}^{2x-4}}{{5}^{2x-3}\ensuremath{\cdot}{3}^{2x-3}}\hfill 
\end{array}
  \end{center}

  
\westep{Subtract exponents (same base)}

\begin{center}
\begin{array}{lcl}
\phantom{\dfrac{{5}^{2x-1}\ensuremath{\cdot}{9}^{x-2}}{{15}^{2x-3}}}& =& {5}^{(2x-1)-(2x-3)}\ensuremath{\cdot}{3}^{(2x-4)-(2x-3)}\hfill \\ 
& =& 5^{2x-1-2x+3} \cdot 3^{2x-4 - 2x+3} \\
& =& {5}^{2}\ensuremath{\cdot}{3}^{-1}\hfill \end{array}
\end{center}


\westep{Skryf die antwoord as a faktor}  
\begin{eqnarray*}
=\frac{25}{3} \\\\ = 8\frac{1}{3}
\end{eqnarray*}

}
\end{wex}

\Note{When working with exponents all the laws of operation for algebra apply.}

\begin{wex}
{%title
Simplifying by taking out a common factor
}
{%question
Vereenvoudig: $$\frac{2^t-2^{t-2}}{3\cdot2^t-2^t} $$

}
{%answer
\westep{Simplify to a form that can be factorised}
\begin{eqnarray*}
\frac{2^t-2^{t-2}}{3\cdot2^t-2^t} & = & \frac{2^t-(2^t\cdot2^{-2})}{3\cdot2^t-2^t} \\
\end{eqnarray*}

\westep{haal uit die gemeenskaplike faktor}

\begin{eqnarray*}
\phantom{\frac{2^t-2^{t-2}}{3\cdot2^t-2^t}}  & = & \frac{2^t(1-2^{-2})}{2^t(3-1)} \\
\end{eqnarray*}


\westep{Cancel common factor and simplify}

\begin{eqnarray*}
\phantom{\frac{2^t-2^{t-2}}{3\cdot2^t-2^t}}  & = & \frac{1-^1/_4}{2} \\
					     & = & \frac{^3/_4}{2} \\
					     & = & \frac{3}{8}
\end{eqnarray*}


} 
\end{wex}




\begin{wex}
{
Vereenvoudig met die gebruik van verskil van kwadrate
}
{
Simplify: 
$$ \frac{9^x-1}{3^x+1} $$
}
{
\westep{Change bases to prime numbers}
\begin{eqnarray*}
 \frac{9^x-1}{3^x+1} & = & \frac{(3^2)^x -1}{3^x+1} \\
		     & = & \frac{(3^x)^2-1}{3^x+1} 
\end{eqnarray*}

\westep{Factorise using difference of squares}

\begin{eqnarray*}
 \phantom{\frac{9^x-1}{3^x+1}} & = & \frac{(3^x-1)(3^x+1)}{3^x+1}\\
\end{eqnarray*}

\westep{Simplfy}

\begin{eqnarray*}
 \phantom{\frac{9^x-1}{3^x+1}} & = & 3^x-1\\
\end{eqnarray*}

}
\end{wex}


\begin{exercises}{Simplifying exponential expressions}
Vereenvoudig sonder om ’n sakrekenaar te gebruik:
\begin{multicols}{2}
\begin{enumerate}[noitemsep, label=\textbf{\arabic*}., itemsep=5pt]
 \item $16^0$
 \item $16a^0$
 \item $\frac{2^{-2}}{3^2}$
 \item $ \frac{5}{2^{-3}}$
 \item $ \Big(\frac{2}{3}\Big)^{-3} $
 \item $ x^2 \cdot x^{3t+1} $
 \item $ 3 \times 3^{2a} \times 3^2$
 \item $ \frac{a^{3x}}{a^x} $
 \item $ \frac{32p^2}{4p^8}$
 \item $ (2t^4)^3$
 \item $ (3^{n+3})^2$
 \item $ \frac{3^n \cdot 9^{n-3}}{27^{n-1}}$
\end{enumerate}
\end{multicols}

 






\end{exercises}




% \par 
% \label{m38359*secfhsst!!!underscore!!!id2193}
% \par \raisebox{-0.2em}{\includegraphics[height=1em]{../icons/www.pdf}} Find the answers with the shortcodes:
%  \par \begin{tabular}[h]{cccccc}
%  (1.) lO6  & \end{tabular}
\section{Rasionaale Eksponente}

We can also apply the exponential laws to expressions with rational exponents.

\begin{wex}
{%title
Vereenvoudig rasionale eksponente
} 
{%question
Simplify: 
$$ (0.008)^{^1/_3} $$
}
{% answer
\westep{Write as a fraction and change bases to prime numbers}

\begin{eqnarray*}
 0.008 & = & \frac{8}{1000} \\
       & = & \frac{2^3}{10^3} \\
       & = & \Big(\frac{2}{10}\Big)^3\\
\end{eqnarray*}

\westep{So}
\begin{eqnarray*}
 (0.008)^{^1/_3} & = & \Big(\Big(\frac{2}{10}\Big)^3\Big)^{^1/_3} \\
		 & = & \frac{2}{10} \\
		 & = & \frac{1}{5}
\end{eqnarray*}
}
\end{wex}

\begin{wex}
{%Title
Vereenvoudig rasionale eksponenete 
}

{
Vereenvoudig:
$$ 2x^{^1/_2} \times 4x^{^{-1}/_{2}} $$
}
{% Answer
\westep{Change the negative exponent}
\begin{eqnarray*}
 2x^{^1/_2} \times 4x^{^{-1}/_{2}} & = & 2x^{^1/_2} \times \frac{4}{x^{^1/_2}} \\
\end{eqnarray*}
Let op dat die $4$ bo bly.\\

\westep{Vereenvoudig}
\begin{eqnarray*}
 \phantom{2x^{^1/_2} \times 4x^{^{-1}/_{2}}} & = & \dfrac{8x^{^1/_2}}{x^{^1/_2}} \\
					     & = & 8
\end{eqnarray*}


}
\end{wex}

\begin{exercises}{Vereenvoudig rasionale eksponenete}
Vereenvoudig sonder om ’n sakrekenaar te gebruik.:
\begin{multicols}{2}
\begin{enumerate}[noitemsep, label=\textbf{\arabic*}., itemsep=5pt]
 \item $ t^{^1/_4} \times 3t^{^7/_4} $
 \item $ \dfrac{16x^2}{(4x^2)^{^1/_2}} $
 \item $ (0.25)^{^1/_2} $
 \item $ (27)^{^{-1}/_3} $
 \item $ (3p^2)^{^1/_2} \times (3p^4)^{^1/_2} $
\end{enumerate}
\end{multicols}

 






\end{exercises}



% \begin{exercises}{Exponential Numbers }
% {
% \label{m38359*id67549}Match the answers to the questions, by filling in the correct answer into the \textbf{Answer} column.
% Possible answers are: $\frac{3}{2}$, $1$, $-1$, $-\frac{1}{3}$, $8$. Answers may be repeated.\par 
% % \textbf{m38359*id67604}\par
%   \begin{table}[H]
% % \begin{table}[H]
% % \\ '' '0'
% \begin{center}
% \label{m38359*id67604}
% \noindent
% \tabletail{%
% \hline
% \multicolumn{2}{|p{\mytableboxwidth}|}{\raggedleft \small \sl continued on next page}\\
% \hline
% }
% \tablelasttail{}
% \begin{xtabular}[t]{|l|l|}\hline
% 	  \textbf{Question}
% 	  &
% 	  \textbf{Answer}
% 	% make-rowspan-placeholders
% \tabularnewline\cline{1-1}\cline{2-2}
% %--------------------------------------------------------------------
% 	  ${2}^{3}$
% 	  &
% % make-rowspan-placeholders
% \tabularnewline\cline{1-1}\cline{2-2}
% %--------------------------------------------------------------------
% 	  ${7}^{3-3}$
% 	  &
% % make-rowspan-placeholders
% \tabularnewline\cline{1-1}\cline{2-2}
% %--------------------------------------------------------------------
% 	  ${\left(\frac{2}{3}\right)}^{-1}$
% 	  &
% % make-rowspan-placeholders
% \tabularnewline\cline{1-1}\cline{2-2}
% %--------------------------------------------------------------------
% 	  ${8}^{7-6}$
% 	  &
% % make-rowspan-placeholders
% \tabularnewline\cline{1-1}\cline{2-2}
% %--------------------------------------------------------------------
% 	  ${\left(-3\right)}^{-1}$
% 	  &
% % make-rowspan-placeholders
% \tabularnewline\cline{1-1}\cline{2-2}
% %--------------------------------------------------------------------
% 	  ${\left(-1\right)}^{23}$
% 	  &
% % make-rowspan-placeholders
% \tabularnewline\cline{1-1}\cline{2-2}
% %--------------------------------------------------------------------
% \end{xtabular}
% \end{center}
% % \begin{center}{\small\bfseries Table 5.1}\end{center}
% % \begin{caption}{\small\bfseries Table 5.1}\end{caption}
% \end{table}
% }
% 
% 
% 
% 
% We will use all these laws in Equations and Inequalities\footnote{\raggedright{}"Equations and Inequalities - Grade 10 [CAPS]" <http://http://cnx.org/content/m38372/latest/>} to help us solve exponential equations.\par \label{m38359*eip-160}

The following video gives an example on using some of the concepts covered in this chapter.
\setcounter{subfigure}{0}
\begin{figure}[H] % horizontal\label{m38359*ExponentsLaw3}
\textnormal{Khan Academy video on Exponents - 5}\vspace{.1in} \nopagebreak
\label{m38359*yt-media5}\label{m38359*yt-video5}
\raisebox{-5 pt}{ \includegraphics[width=0.5cm]{col11306.imgs/summary_www.png}} { (Video:  MG10049 )}
% \vspace{2pt}
% \vspace{.1in}
\end{figure}    


\section{Exponential Equations}

Exponential equations have the unknown variable in the exponent. Here are some examples:
\begin{eqnarray*}
 3^{x+1} & = & 9 \\
5^t + 3 \cdot 5^{t-1} & = & 400
\end{eqnarray*}

Solving exponential equations is simple, we need to apply the laws of exponents.

\Note{Indien $a>0$ en $a \ne 0$ en
$$ a^x = a^y $$
dan
$$ x = y ~~\mbox{(same base)}$$
(Indien $a=1$, dan kan $x$ en $y$ verskil)
}


This means that if we can write a single term with the same base on each side of the equation, we can equate the exponents.


\begin{wex}
{%title
Equating exponents
}
{%question
Solve $x$
$$ 3^{x+1} = 9 $$
}
{%answer

\westep{Change bases prime numbers}
\begin{eqnarray*}
 3^{x+1} & = & 3^2 \\
\end{eqnarray*}

\westep{Bases are the same so we can equate exponents}
\begin{eqnarray*}
 {x+1} & = & 2 \\
\therefore x & = & 1
\end{eqnarray*}
}
\end{wex}

\Note{To solve exponential equations, we use all the strategies for solving linear and quadratic equations.}



\begin{wex}
{%title
Solving equations by taking out a common factor
}
{%question
Solve $t$
$$ 5^t + 3 \cdot 5^{t+1} = 400 $$
}
{%answer

\westep{Rewrite the expression}
\begin{eqnarray*}
 5^t + 3 ( 5^t \cdot 5) & = & 400 \\
\end{eqnarray*}

\westep{Take out a common factor}
\begin{eqnarray*}
 5^t(1 + 3 \cdot 5) & = & 400 \\
\end{eqnarray*}


\westep{Vereenvoudig}

\begin{eqnarray*}
 5^t(16) & = & 400 \\
  5^t & = & 25 \\
\end{eqnarray*}


\westep{Change bases prime numbers}
\begin{eqnarray*}
  5^t & = & 5^2 \\
\end{eqnarray*}


\westep{Bases are the same so we can equate exponents}
\begin{eqnarray*}
\therefore t & = & 2
\end{eqnarray*}

}
\end{wex}

\begin{wex}
{%title
Solving equations by factorising a trinomial
}
{%question
Solve
$$ p-13 p^{^1/_2} + 36 =  0$$
}
{ % Answer

\westep{We notice that $(p^{^1/_2})^2=p$ so we can rewrite the equation as}

$$ (p^{^1/_2})^2 -13p^{^1/_2} + 36 = 0 $$

\westep{Factorise as a trinomial}

$$ (p^{^1/_2} -9)(p^{^1/_2}-4) = 0 $$

\westep{Solve to find both roots}

%Left Column				%right column
\begin{align*}
p^{^1/_2} - 9 &= 0			&   p^{^1/_2} - 4 &= 0		\\
p^{^1/_2} &= 9				&   p^{^1/_2} &= 4		\\		
(p^{^1/_2})^2 &= (9)^2			&   (p^{^1/_2})^2 &= (4)^2\\
p &= 9^2				&   p &= 16\\
p &= 81					&\\
\end{align*} 
Therefore $p=81$ or $p=16$


}
\end{wex}

\begin{exercises}{Solving exponential equations}
Solve:
\begin{multicols}{2}
\begin{enumerate}[noitemsep, label=\textbf{\arabic*}., itemsep=5pt]
\item $ 2^{x+5} = 32 $
\item $ 5^{2x+2} = \dfrac{1}{125} $
\item $ 64^{y+1} = 16^{2y+5} $
\item $ 3^{9x-2} = 27 $
\item $ 81^{k+2} = 27^{k+4} $
\setcounter{enumi}{6}
\item $ 25^{(1-2x)}-5^4 = 0 $
\item $ 27^x \times 9^{x-2} = 1 $
\item $ 2^t + 2^{t+2} = 40 $
\item $ 2 \cdot 5^{2-x} = 5+ 5^X $
\item $ 9^m + 3^{3-2m} = 28 $
\item $ y - 2y^{^1/_2} + 1 = 0 $
\end{enumerate}
\end{multicols}

\begin{enumerate}[noitemsep, label=\textbf{\arabic*}., itemsep=5pt]
\setcounter{enumi}{5}
 \item The growth of an algae can be modelled by the function $f(t) = 2^t$. Find the value of $t$ such that $f(t)=128$.   
\end{enumerate}
\end{exercises}




\section{ Opsomming}
\label{m38359*eip-908} $ \hspace{-5pt}\begin{array}{cccccccccccc}   \end{array} $ \hspace{2 pt}\raisebox{-0.2em}{\includegraphics[height=1em]{../icons/www.pdf}} {(section shortcode: MG10050 )} 


\begin{itemize}[noitemsep, label=\textbullet{}]
    \item Eksponensiaalnotasie verwys na ’n getal wat geskryf word as ${a}^{n}$ waar $n$ ’n heelgetal is en $a$ enige reële getal is.
    \item $a$ is die \textsl{grondtal} en $n$ is die \textsl{eksponent} or \textsl{index}.
    \item Definisie: 
	  \begin{itemize}[noitemsep]
	   \item ${a}^{n}=a\ensuremath{\times}a\ensuremath{\times}\cdots \ensuremath{\times}a\phantom{\rule{2.em}{0ex}}\left(\mbox{$n$ times}\right)$
	   \item ${a}^{0}=1,a\ne 0$
	   \item ${a}^{-n}=\frac{1}{{a}^{n}},a\ne 0$
	  \end{itemize}

    
    \item  Die Eksponentewette: 
	\begin{itemize}[itemsep=4pt]
	    \item  ${a}^{m}\ensuremath{\times}{a}^{n}={a}^{m+n}$
	    \item  ${\dfrac{{a}^{m}}{{a}^{n}}={a}^{m-n}$
	    \item  ${\left(ab\right)}^{n}={a}^{n}{b}^{n}$
            \item  $\Big(\dfrac{a}{b}\Big)^n = \dfrac{a^n}{b^n}$
	    \item  ${\left({a}^{m}\right)}^{n}={a}^{mn}$
	\end{itemize}
    \end{itemize}


\begin{eocexercises}{}
%  $ \hspace{-5pt}\begin{array}{cccccccccccc}   \end{array} $ 
% \hspace{2 pt}\raisebox{-0.2em}{\includegraphics[height=1em]{../icons/www.pdf}} {(section shortcode: MG10051 )} \par 
% \label{m38359*id67892}


Vereenvoudig:
\begin{multicols}{2}
\begin{enumerate}[noitemsep, label=\textbf{\arabic*}., itemsep=5pt]
\item $ t^3 \times 2t^0 $
\item $ 5^{2x+y} \cdot 5^{3(x+z)} $
\item $ (b^{k+1})^k $
\item $ \dfrac{6^{5p}}{9^p} $
\item $ m^{-2t} \times (3m^t)^3 $
\item $\dfrac{3{x}^{-3}}{{\left(3x\right)}^{2}}$
\item $\dfrac{{5}^{b-3}}{{5}^{b+1}}$
\item $\dfrac{{2}^{a-2}.{3}^{a+3}}{{6}^{a}}$
\item $\dfrac{{3}^{n}\ensuremath{\cdot}{9}^{n-3}}{{27}^{n-1}}$
\item ${\left(\dfrac{2{x}^{2a}}{{y}^{-b}}\right)}^{3}$
\item $\dfrac{{2}^{3x-1}\ensuremath{\cdot}{8}^{x+1}}{{4}^{2x-2}}$
\item $\dfrac{{6}^{2x}\ensuremath{\cdot}{11}^{2x}}{{22}^{2x-1}\ensuremath{\cdot}{3}^{2x}}$
\item $\dfrac{{\left(-3\right)}^{-3}\ensuremath{\cdot}{\left(-3\right)}^{2}}{{\left(-3\right)}^{-4}}$
\item ${\left({3}^{-1}+{2}^{-1}\right)}^{-1}$
\item $\dfrac{{9}^{n-1}\ensuremath{\cdot}{27}^{3-2n}}{{81}^{2-n}}$
\item $\dfrac{{2}^{3n+2}\ensuremath{\cdot}{8}^{n-3}}{{4}^{3n-2}}$
\item $\dfrac{3^{t+3} + 3^t}{2 \cdot 3^t} $
\item $\dfrac{2^{3p} +1}{2^p + 1} $
\end{enumerate}
\end{multicols}



Solve:
\begin{multicols}{2}
\begin{enumerate}[noitemsep, label=\textbf{\arabic*}., itemsep=5pt]
\setcounter{enumi}{18}
\item $ 3^x = \dfrac{1}{27} $
\item $ 5^{t-1} = 1 $
\item $ 2 \cdot 7^{3x} = 98 $
\item $ 2^{m+1} = (0.5)^{m-2}$
\item $ 3^{y+1} = 5^{y+1} $
\item $ z^{^3/_2} = 64 $
\item $ 16x^{^1/_2} - 4 = 0 $
\item $ m^0 + m^{-1} = 0 $
\item $ t^{^1/_2} - 3t^{^1/_4} + 2 = 0 $
\item $ 3^p + 3^p + 3^p = 27 $
\item $ k^{-1} - 7x^{^{-1}/_2} -18 = 0 $
\item $ x^{^1/_2}+3x^{^1/_4}-18 = 0 $
\end{enumerate}
\end{multicols}


\end{eocexercises}

% \begin{enumerate}[noitemsep, label=\textbf{\arabic*}. ] 
% %     \item Simplify as far as possible:
% % 	\begin{enumerate}[noitemsep, label=\textbf{\alph*}. ] 
% % 	    \item ${302}^{0}$
% % 	    \item ${1}^{0}$
% % 	    \item ${\left(xyz\right)}^{0}$
% % 	    \item ${\left[{\left(3{x}^{4}{y}^{7}{z}^{12}\right)}^{5}{\left(-5{x}^{9}{y}^{3}{z}^{4}\right)}^{2}\right]}^{0}$
% % 	    \item ${\left(2x\right)}^{3}$
% % 	    \item ${\left(-2x\right)}^{3}$
% % 	    \item ${\left(2x\right)}^{4}$
% % 	    \item ${\left(-2x\right)}^{4}$
% % 	\end{enumerate}
%     \item Simplify
%     \item Simplify without using a calculator. Leave your answers with positive exponents.
% 	\begin{enumerate}[noitemsep, label=\textbf{\alph*}. ] 
% 	    \item $\dfrac{3{x}^{-3}}{{\left(3x\right)}^{2}}$
% 	    \item $5{x}^{0}+{8}^{-2}-{\left(\frac{1}{2}\right)}^{-2}\ensuremath{\cdot}{1}^{x}$
% 	    \item $\dfrac{{5}^{b-3}}{{5}^{b+1}}$
% 	\end{enumerate}
%     \item Simplify, showing all steps:
% 	\begin{enumerate}[noitemsep, label=\textbf{\alph*}. ] 
% 	    \item $\dfrac{{2}^{a-2}.{3}^{a+3}}{{6}^{a}}$
% 	    \item $\dfrac{{a}^{2m+n+p}}{{a}^{m+n+p}\ensuremath{\cdot}{a}^{m}}$
% 	    \item $\dfrac{{3}^{n}\ensuremath{\cdot}{9}^{n-3}}{{27}^{n-1}}$
% 	    \item ${\left(\dfrac{2{x}^{2a}}{{y}^{-b}}\right)}^{3}$
% 	    \item $\dfrac{{2}^{3x-1}\ensuremath{\cdot}{8}^{x+1}}{{4}^{2x-2}}$
% 	    \item $\dfrac{{6}^{2x}\ensuremath{\cdot}{11}^{2x}}{{22}^{2x-1}\ensuremath{\cdot}{3}^{2x}}$
% 	\end{enumerate}
%     \item Simplify, without using a calculator:
% 	\begin{enumerate}[noitemsep, label=\textbf{\alph*}. ] 
% 	    \item $\dfrac{{\left(-3\right)}^{-3}\ensuremath{\cdot}{\left(-3\right)}^{2}}{{\left(-3\right)}^{-4}}$
% 	    \item ${\left({3}^{-1}+{2}^{-1}\right)}^{-1}$
% 	    \item $\dfrac{{9}^{n-1}\ensuremath{\cdot}{27}^{3-2n}}{{81}^{2-n}}$
% 	    \item $\dfrac{{2}^{3n+2}\ensuremath{\cdot}{8}^{n-3}}{{4}^{3n-2}}$
% 	\end{enumerate}
% \end{enumerate}
% 
% \raisebox{-0.2em}{\includegraphics[height=1em]{../icons/www.pdf}} Find the answers with the shortcodes:
% \begin{tabular}[h]{cccccc}
% (1.) lOJ  &  (2.) lOu  &  (3.) lOS  &  (4.) lOh  & 
% \end{tabular}
